% =========================================================
% The Logarithm --- Notes
% Chapter: elementary-functions
% Section: logarithmic
% Volume III $\cdot$ Analysis
% =========================================================

\subsection*{The Logarithm}
\label{sec:logarithmic}

% ---------------------------------------------------------
% Toolkit
% ---------------------------------------------------------
% ---------------------------------------------------------
% Definition --- Natural Logarithm
% ---------------------------------------------------------
\begin{definitionbox}{Definition (Natural Logarithm)}
\begin{definition}[Natural Logarithm]
\label{def:natural-logarithm}
The \textbf{natural logarithm} $\ln : (0,\infty) \to \mathbb{R}$ is
the inverse of the exponential function: for $x > 0$,
\[
  \ln x \;:=\; \text{the unique } y \in \mathbb{R}
  \text{ such that } e^y = x.
\]
Equivalently, $\ln x$ is defined by the integral
\[
  \ln x \;:=\; \int_1^x \frac{1}{t}\, dt.
\]
\end{definition}
\end{definitionbox}

\begin{remark*}[Standard quantified statement]
\[
  \forall x > 0 :\; e^{\ln x} = x.
  \qquad
  \forall y \in \mathbb{R} :\; \ln(e^y) = y.
\]
\end{remark*}

\begin{remark*}[Predicate reading]
\[
  y=\ln x
  \;\colonequiv\;
  x > 0 \;\wedge\; e^y = x
  \;\colonequiv\;
  x > 0 \;\wedge\; y = \int_1^x \frac{1}{t}\,dt.
\]


\[
  \operatorname{TendsTo}\!\left(\frac{\ln(1+x)}{x},\,0,\,1,\,\mathbb R,\,\mathbb R\right)
  \;\colonequiv\;
  \forall\varepsilon>0,\;\exists\delta>0,\;\forall x:\;
  \bigl(
    \operatorname{IsPuncturedCenteredOpenInterval}(I^\times_\delta(0),I_\delta(0),0,\delta,\mathbb R)
    \wedge
    \operatorname{Within}(x,I^\times_\delta(0),\mathbb R)
  \bigr)
  \Rightarrow
  \bigl(
    \operatorname{IsCenteredOpenInterval}(I_\varepsilon(1),1,\varepsilon,\mathbb R)
    \wedge
    \operatorname{Within}\!\left(\frac{\ln(1+x)}{x},I_\varepsilon(1),\mathbb R\right)
  \bigr).
\]
\end{remark*}

\begin{remark*}[Negated quantified statement]
\[
  \neg(y=\ln x)
  \;\iff\;
  x \leq 0 \;\vee\; (x>0 \wedge y\neq\ln x).
\]
\end{remark*}

\begin{remark*}[Negation predicate reading]
\[
\neg(y=\ln x)
  \;\iff\;
  x \leq 0 \;\vee\; (x>0 \wedge y\neq\ln x).
\]
\end{remark*}

\begin{remark*}[Failure modes]
\begin{description}
  \item[Exposition.]
  The negated statement isolates the obstruction to the asserted condition.

  \item[\neq-condition.]
  The displayed obstruction is the corresponding failure mode.
  \[
\neg(y=\ln x)
  \;\iff\;
  x \leq 0 \;\vee\; (x>0 \wedge y\neq\ln x).
  \]
  \[
\neg(y=\ln x)
  \;\iff\;
  x \leq 0 \;\vee\; (x>0 \wedge y\neq\ln x).
  \]
\end{description}
\end{remark*}



\begin{remark*}[Interpretation]
The inverse-of-exponential definition requires $e^x$ to be
established first. The integral definition is self-contained and
is preferred when building analysis from scratch. Both yield the
same function.


\begin{enumerate}[label=(\roman*)]
  \item $\ln 1 = 0$;\; $\ln e = 1$.
  \item \textbf{Functional equation:} $\ln(xy) = \ln x + \ln y$
    for all $x, y > 0$.
  \item $\ln(x/y) = \ln x - \ln y$;\;
    $\ln(x^r) = r\ln x$ for $r \in \mathbb{R}$.
  \item $\ln$ is strictly increasing on $(0,\infty)$.
  \item $\ln x \to +\infty$ as $x \to \infty$;\;
    $\ln x \to -\infty$ as $x \to 0^+$.
  \item $(\ln x)' = 1/x$ for all $x > 0$.
\end{enumerate}


For $a > 0$, $a \neq 1$, the \textbf{logarithm base $a$} is
$\log_a x := \ln x / \ln a$. The common logarithm is $\log_{10}$;
the binary logarithm is $\log_2$.
\end{remark*}

\begin{dependencies}
\begin{itemize}
  \item \hyperref[def:exponential]{Exponential Function}
  \item \hyperref[def:riemann-integral]{Riemann Integral}
\end{itemize}
\end{dependencies}

% ---------------------------------------------------------
% Proposition --- Fundamental Logarithmic Limit
% ---------------------------------------------------------

\begin{propositionbox}{Proposition (Fundamental Logarithmic Limit)}
\begin{proposition}[Fundamental Logarithmic Limit]
\label{prop:limit-logarithm-fundamental}
\[
  \lim_{x \to 0} \frac{\ln(1+x)}{x} = 1.
\]

\smallskip\noindent


\hyperref[prf:limit-logarithm-fundamental]{\textit{Go to proof.}}
\end{proposition}
\end{propositionbox}

\begin{remark*}[Standard quantified statement]
\[
  \forall\varepsilon > 0,\;
  \exists\delta > 0 :\;
  0 < |x| < \delta \Rightarrow
  \left|\frac{\ln(1+x)}{x} - 1\right| < \varepsilon.
\]
\end{remark*}

\begin{remark*}[Negated quantified statement]
\[
  \neg\left(\lim_{x\to 0}\frac{\ln(1+x)}{x}=1\right)
  \;\iff\;
  \exists\varepsilon_0>0,\;\forall\delta>0,\;\exists x:\;
  0<|x|<\delta \;\wedge\;
  \left|\frac{\ln(1+x)}{x}-1\right|\geq\varepsilon_0.
\]
\end{remark*}

\begin{remark*}[Negation predicate reading]
\[
\neg\left(\lim_{x\to 0}\frac{\ln(1+x)}{x}=1\right)
  \;\iff\;
  \exists\varepsilon_0>0,\;\forall\delta>0,\;\exists x:\;
  0<|x|<\delta \;\wedge\;
  \left|\frac{\ln(1+x)}{x}-1\right|\geq\varepsilon_0.
\]
\end{remark*}

\begin{remark*}[Failure modes]
\begin{description}
  \item[Exposition.]
  The negated statement isolates the obstruction to the asserted condition.

  \item[\(\exists\)-condition.]
  The displayed obstruction is the corresponding failure mode.
  \[
\neg\left(\lim_{x\to 0}\frac{\ln(1+x)}{x}=1\right)
  \;\iff\;
  \exists\varepsilon_0>0,\;\forall\delta>0,\;\exists x:\;
  0<|x|<\delta \;\wedge\;
  \left|\frac{\ln(1+x)}{x}-1\right|\geq\varepsilon_0.
  \]
  \[
\neg\left(\lim_{x\to 0}\frac{\ln(1+x)}{x}=1\right)
  \;\iff\;
  \exists\varepsilon_0>0,\;\forall\delta>0,\;\exists x:\;
  0<|x|<\delta \;\wedge\;
  \left|\frac{\ln(1+x)}{x}-1\right|\geq\varepsilon_0.
  \]
\end{description}
\end{remark*}



\begin{remark*}[Interpretation]
The fundamental exponential limit $\lim(e^x-1)/x = 1$ and the
fundamental logarithmic limit $\lim\ln(1+x)/x = 1$ are equivalent:
each follows from the other via the substitution $x \leftrightarrow
e^x - 1$. Both express the first-order Taylor approximation of the
respective function at $0$.


From the series: $\ln(1+x) = x - x^2/2 + x^3/3 - \cdots$ (valid
for $|x| < 1$). Dividing by $x$ gives $1 - x/2 + x^2/3 - \cdots
\to 1$ as $x \to 0$ by continuity of power series.
\end{remark*}

\begin{dependencies}
\begin{itemize}
  \item \hyperref[def:exponential]{Exponential Function}
  \item \hyperref[prop:inverse-bijection]{Inverse Bijection}
  \item \hyperref[thm:continuous-inverse-theorem]{Continuous Inverse Theorem}
\end{itemize}
\end{dependencies}

% ---------------------------------------------------------
% Proposition --- Every Power Dominates the Logarithm
% ---------------------------------------------------------

\begin{propositionbox}{Proposition (Every Power Dominates the Logarithm)}
\begin{proposition}[Every Power Dominates the Logarithm]
\label{prop:power-dominates-log}
For every $r > 0$:
\[
  \lim_{x \to \infty} \frac{\ln x}{x^r} = 0
  \qquad\text{and}\qquad
  \lim_{x \to 0^+} x^r \ln x = 0.
\]

\smallskip\noindent


\hyperref[prf:power-dominates-log]{\textit{Go to proof.}}
\end{proposition}
\end{propositionbox}

\begin{remark*}[Standard quantified statement]
\[
  \forall r > 0 :\;
  \lim_{x \to \infty} x^{-r} \ln x = 0
  \;\wedge\;
  \lim_{x \to 0^+} x^r \ln x = 0.
\]
\end{remark*}

\begin{remark*}[Predicate reading]
\[
  \operatorname{LimitAtInfinity}(x^{-r}\ln x,\,+\infty,\,0)
  \;\colonequiv\;
  \forall\varepsilon>0,\;\exists M\in\mathbb{R},\;\forall x:\;
  x>M \Rightarrow
  \bigl(
    \operatorname{IsCenteredOpenInterval}(I_\varepsilon(0),0,\varepsilon,\mathbb R)
    \wedge
    \operatorname{Within}(x^{-r}\ln x,I_\varepsilon(0),\mathbb R)
  \bigr).
\]
\end{remark*}

\begin{remark*}[Negated quantified statement]
\[
  \neg\left(\lim_{x\to\infty}\frac{\ln x}{x^r}=0\right)
  \;\iff\;
  \exists\varepsilon_0>0,\;\forall M>0,\;\exists x>M:\;
  x^{-r}|\ln x|\geq\varepsilon_0.
\]
\end{remark*}

\begin{remark*}[Negation predicate reading]
\[
\neg\left(\lim_{x\to\infty}\frac{\ln x}{x^r}=0\right)
  \;\iff\;
  \exists\varepsilon_0>0,\;\forall M>0,\;\exists x>M:\;
  x^{-r}|\ln x|\geq\varepsilon_0.
\]
\end{remark*}

\begin{remark*}[Failure modes]
\begin{description}
  \item[Exposition.]
  The negated statement isolates the obstruction to the asserted condition.

  \item[\(\exists\)-condition.]
  The displayed obstruction is the corresponding failure mode.
  \[
\neg\left(\lim_{x\to\infty}\frac{\ln x}{x^r}=0\right)
  \;\iff\;
  \exists\varepsilon_0>0,\;\forall M>0,\;\exists x>M:\;
  x^{-r}|\ln x|\geq\varepsilon_0.
  \]
  \[
\neg\left(\lim_{x\to\infty}\frac{\ln x}{x^r}=0\right)
  \;\iff\;
  \exists\varepsilon_0>0,\;\forall M>0,\;\exists x>M:\;
  x^{-r}|\ln x|\geq\varepsilon_0.
  \]
\end{description}
\end{remark*}



\begin{remark*}[Interpretation]
For the first: substitute $x = e^t$ to get $t/e^{rt} \to 0$ by
exponential dominance. For the second: substitute $x = 1/t$ and
reduce to the first: $x^r \ln x = (-\ln t)/t^r \to 0$.


These limits establish the position of $\ln x$ in the growth
hierarchy: $\ln x$ is dominated by every positive power. Combined
with the exponential dominance result, the full ordering is
$\ln x \ll x^r \ll e^x$ as $x \to \infty$ for any $r > 0$.
The second limit says that near $0^+$, the positive power $x^r$
wins over the logarithmic blow-up $|\ln x| \to \infty$. Used in
integrability arguments near $0$.
\end{remark*}

\begin{dependencies}
\begin{itemize}
  \item \hyperref[def:natural-logarithm]{Natural Logarithm}
  \item \hyperref[def:power-function]{Power Function}
  \item \hyperref[def:limit-function]{Function Limit}
  \item \hyperref[def:limit-at-infinity]{Limit at Infinity}
  \item \hyperref[prop:limit-exponential-fundamental]{Fundamental Exponential Limit}
  \item \hyperref[prop:exponential-dominates-power]{Exponential Dominates Every Power}
\end{itemize}
\end{dependencies}
